\clearpage
\section{RX-side Conversion Gain --- High-Side LO (Spectral Flip)}
\label{sec:rx-gain-highside}

\textbf{What is conversion gain?} Exactly as in Section~\ref{sec:rx-gain-lowside}: the power of the signal
coming out of the ``RX-side FR1 IF Out'' connector, minus the power of the signal fed into the ``RX-side FR3 RF In'' connector, both in dBm,
gives the conversion gain in dB. Only the LO's position changes here, now above FR3 instead of below it.

\textbf{What's plotted.} Identical methodology to Section~\ref{sec:rx-gain-lowside}, with the unit tuned to
high-side LO instead: each of the first seven figures holds FR1 (the IF output frequency) fixed at one center
value (1.0 up to the unit's 4.0\,GHz maximum, in 0.5\,GHz steps) and plots conversion gain against the swept
FR3 input frequency for that output, using the same methodology as the TX conversion-gain sections: every point
is the average conversion gain across a real 500\,MHz swept window centered on that point, not a single CW
tone.\footnote{Every point is the average conversion gain across a real 500\,MHz swept window, at a calibrated
−50\,dBm input level, same as Section~\ref{sec:rx-gain-lowside}.} Grey is the raw LNA-and-mixer chain with
the AM3153 filter bypassed; blue is the filter engaged, re-selected at every point using the same ``optimal''
algorithm described in Section~\ref{sec:sideband-suppression}, reused here because the AM3153 is RX's actual
image-rejection mechanism --- there is no I/Q image rejection on this side of the board, since the mixer's Q
output is simply discarded. The gap between the two traces, at any point, is exactly how much the filter costs
in insertion loss there, same as on the TX-side gain sections.

\textbf{Measurement conditions.} The same as in Section~\ref{sec:rx-gain-lowside} (the same run): LMX2820 die at \condRxGainMinC--\condRxGainMaxC\textdegree C.

\textbf{How the filter is chosen.} The same two-step rule used throughout this datasheet: first find whichever
of the 48 candidate configurations (3 bands × 16 sub-states) passes the most desired energy on its own;
then, restricted to candidates within 3\,dB of that best desired power, choose whichever rejects the image the
most. Consecutive points that land on the same filter form a \emph{region}; the dashed vertical lines on each
panel mark every boundary between regions, and the blue hex code running up from the x-axis inside each region
is the actual \texttt{filter\_bits} value used for every point in it --- so, exactly as on the TX side, the
exact filter behind any measurement on these panels is always readable directly off the figure, with no need to
consult a separate table.

\textbf{High-side LO: LO = FR3 + FR1.} This is the LO tuning regime shown in this section. Because of it,
the IF output spectrum \emph{is} flipped (mirrored) relative to the RF input --- as FR3 increases, the IF
output's own internal content decreases.\footnote{Since TX and RX share one physical LO, whenever RX is
running high-side LO, TX (if active at the same time) is necessarily running lower-sideband up-conversion, and
vice versa.} Every point's window is exactly 500\,MHz wide, the same fixed window used throughout the TX-side
gain sections: a fixed filter setting's measured response across that full window is exactly what a real
wideband signal at that setting would experience. Figure \ref{fig:8-7} (FR1 = 4.0\,GHz) is the one exception,
just as on the TX side: its window extends 250\,MHz past the unit's IF design spec, called out directly on
that panel and outlined in red on Figures \ref{fig:8-8} and \ref{fig:8-9}, which repeat the bypassed and
optimal data across the full measured (FR1, FR3) grid for a single-glance view instead of the seven separate
slices above.

\figpair{images/sec8_fig1.pdf}{FR1 (output) = 1.0\,GHz.}{fig:8-1}%
        {images/sec8_fig2.pdf}{FR1 (output) = 1.5\,GHz.}{fig:8-2}

\figpair{images/sec8_fig3.pdf}{FR1 (output) = 2.0\,GHz.}{fig:8-3}%
        {images/sec8_fig4.pdf}{FR1 (output) = 2.5\,GHz.}{fig:8-4}
\figpair{images/sec8_fig5.pdf}{FR1 (output) = 3.0\,GHz.}{fig:8-5}%
        {images/sec8_fig6.pdf}{FR1 (output) = 3.5\,GHz.}{fig:8-6}
\figonesmall{images/sec8_fig7.pdf}{FR1 (output) = 4.0\,GHz. 250\,MHz of this window falls outside the unit's IF design spec (shown in red).}{fig:8-7}

\newpage
\figrow{images/sec8_fig8.pdf}{Conversion gain, high-side LO: full measured (FR1, FR3) grid, filter bypassed.
Zoom in to read each cell's exact values, printed as (FR1, FR3, conversion gain), with FR1 and FR3 in GHz and
the conversion gain in dB.}{fig:8-8}
\figrow{images/sec8_fig9.pdf}{Conversion gain, high-side LO: full measured (FR1, FR3) grid, optimal filter.
Zoom in to read each cell's exact values, printed as (FR1, FR3, conversion gain), with FR1 and FR3 in GHz and
the conversion gain in dB.}{fig:8-9}
